% ─── Chapter 8: Results & Analysis ─────────────────────────────────
\chapter{Results \& Analysis}
\label{ch:results}

This chapter synthesises the experimental outcomes reported in
Chapter~\ref{ch:experiments} into an overall assessment of training performance.

\section{Curriculum Stage Results}

\subsection{Stage 1: Random Warmup}

The agent reliably solves the easy stage---achieving the promotion threshold
(70\,\% win rate) against the \texttt{RandomNoSwitchPlayer} opponent---within
approximately 84\,000~training steps.  This confirms that the training
pipeline, observation encoding, and basic model architecture function
correctly.

% TODO: Insert MLflow win rate curve for easy stage showing rapid convergence
% to >70% win rate.  Mark the promotion point.

\subsection{Stage 2: Self-Play}

% TODO: Update this section with self-play training results once available.
% Include: win rate progression during self-play stage, comparison with
% non-self-play baseline, any observed oscillation or instability in
% self-play dynamics, and whether self-play helped break the plateau.
% Placeholder talking points:
% - Self-play introduces a natural curriculum where difficulty increases
%   as the agent improves.
% - The SelfPlayPlayer hot-reloads weights from the latest checkpoint,
%   so opponents are always close to the agent's current skill level.
% - Potential issue: oscillation if the agent learns to exploit its own
%   weaknesses, then those weaknesses disappear in the next update.

\subsection{Stage 3: Mixed Opponents}

Against the mixed opponent pool (50\,\% heuristic, 30\,\% self-play, 20\,\%
random), the agent reaches approximately 50\,\% overall win rate but does not
improve beyond this level even after 2.5~million steps.  The win rate breaks
down as:

\begin{center}
\small
\begin{tabularx}{0.7\textwidth}{l c}
  \toprule
  \textbf{Opponent} & \textbf{Win Rate} \\
  \midrule
  Random      & $\sim$95\,\% \\
  Heuristic   & $\sim$5\,\% \\
  Self-play   & Varies (see above) \\
  \bottomrule
\end{tabularx}
\end{center}

The 50\,\% overall rate is an artefact of the opponent mixture: the agent wins
almost all random-opponent games and loses almost all heuristic-opponent games.

% TODO: Insert MLflow chart showing win rate separated by opponent type over
% training steps.  The random curve should be high and stable; the heuristic
% curve should be near zero throughout.

\section{Validation Checkpoint Results}

Fixed-paired validation at key checkpoints:

\begin{center}
\small
\begin{tabularx}{\textwidth}{r c c c}
  \toprule
  \textbf{Checkpoint} & \textbf{Overall} & \textbf{Vs.\ Random} & \textbf{Vs.\ Heuristic} \\
  \midrule
  Step 51\,200   & 37.5\,\% & 75\,\%  & 0\,\%  \\
  Step 200K      & 45\,\%   & 90\,\%  & 0\,\%  \\
  Final (2.5M)   & 50\,\%   & 95\,\%  & 5\,\%  \\
  \bottomrule
\end{tabularx}
\end{center}

% TODO: Insert MLflow validation comparison chart across checkpoints.
% Show per-opponent-type bars or lines.

The progression shows consistent improvement against random opponents but
negligible progress against heuristic ones, confirming that the learned
strategy does not generalise to more competent play.

\section{Loss Metrics Analysis}

A consistent loss signature has been observed across multiple training runs
during the medium/mixed stage:

\begin{itemize}
  \item \textbf{Policy loss:} Near zero and flat.  The policy network is not
        making meaningful gradient updates---it has converged to a local
        optimum that beats random opponents but not heuristic ones.
  \item \textbf{Value loss:} High and noisy.  The value network cannot
        accurately predict expected returns, resulting in poor GAE advantage
        estimates.
  \item \textbf{Entropy:} Remains high ($>$1.5 nats).  The policy does not
        converge to a confident action distribution, suggesting it has not
        identified a reliable strategy.
  \item \textbf{Explained variance:} Negative ($-0.003$).  The value
        function's predictions are worse than simply outputting the mean
        return.
\end{itemize}

% TODO: Insert MLflow loss curves (policy loss, value loss, entropy)
% over the full training run.  Annotate with curriculum stage transitions.

\section{Attention Analysis}

After the architectural improvements described in Chapter~\ref{ch:experiments},
the final single-layer transformer with cross-team attention bias produces
healthy attention patterns:

\begin{center}
\small
\begin{tabularx}{0.7\textwidth}{l r}
  \toprule
  \textbf{Attention Pair} & \textbf{Weight} \\
  \midrule
  CLS $\to$ Opponent active     & 37.5--62.5\,\% \\
  Our active $\to$ Opponent active & 48.9\,\% \\
  \bottomrule
\end{tabularx}
\end{center}

This is a significant improvement over the 4-layer baseline, where opponent
attention collapsed to 3\,\% by layer~1.  The model is now clearly
``looking at'' the opponent.  However, this improved attention has not
translated into improved win rates against heuristic opponents.

% TODO: Insert attention heatmap comparison:
% Left: 4-layer baseline (collapsed attention)
% Right: 1L + attention bias (structured, opponent-aware attention)

\section{Key Insight: Architecture vs.\ Reward Signal}

The central finding from the experimental programme is:

\begin{keyinsight}
\textbf{Architecture is no longer the bottleneck.}  The model can see the
opponent, distinguish tokens by position and role, use stable vocabulary
embeddings, and operate in a compressed action space.  The remaining obstacle
is the \textbf{reward signal}: the current reward components do not
sufficiently incentivise the agent to develop strategic behaviour against
competent opponents.  The value function fails to learn (negative explained
variance), preventing policy improvement through PPO's advantage estimation.
\end{keyinsight}

This redirects the optimisation focus from architecture tuning to reward
design, potentially including:
\begin{itemize}
  \item Auxiliary tasks (e.g.\ predicting opponent moves, predicting battle
        outcome from intermediate states).
  \item More granular shaped rewards that capture tactical elements (e.g.\
        setting up favourable type matchups, preserving key Pok\'{e}mon).
  \item Reward scaling or normalisation to improve value-function learning.
\end{itemize}

\section{Honest Assessment}

At the time of writing, the agent demonstrates:

\begin{center}
\small
\begin{tabularx}{\textwidth}{l X}
  \toprule
  \textbf{What works} & \textbf{What does not} \\
  \midrule
  Training pipeline runs end-to-end &
    Agent cannot beat heuristic opponents consistently \\
  Easy stage solved reliably &
    Value function fails to learn (negative explained variance) \\
  Observation bugs found and fixed &
    Reward signal does not incentivise strategic play \\
  Attention collapse diagnosed and resolved &
    No automated gauntlet validation yet \\
  Distributed infrastructure stable at 48 envs &
    LSTM not validated end-to-end in training \\
  MLflow tracking produces complete metrics &
    Milestone~1 competency not demonstrated \\
  \bottomrule
\end{tabularx}
\end{center}

The engineering foundation is solid.  The remaining challenge is a core ML
problem: designing a reward function that produces a learnable value landscape
for a complex, sparse-reward environment.
